\chapter{Introduction}
\label{chap:Introduction}

\section{Motivation}
\label{sec:Motivation}

Large language models have become competent at agentic tasks, in which the model
plans over multiple turns, issues calls to external tools, and conditions its
subsequent behaviour on the results those tools return. The models that perform best
in this setting are large, expensive to serve, and in most cases proprietary. A
practitioner who needs an agent to run at low latency, at high request volume, on
private infrastructure, or under a fixed inference budget cannot in general deploy
the model that performs the task best. The practical question is therefore how to
transfer agentic competence from a large teacher model into a small student model
that can be deployed under those constraints.

Supervised fine-tuning on teacher demonstrations is the standard answer, and it is
insufficient for agentic tasks in a specific and well-understood way. A student
trained on the teacher's trajectories is optimized only on the states the teacher
visits. At inference time the student visits its own states, and any deviation
compounds: an agent that issues a malformed tool call in its third turn spends the
remainder of the episode in a region of the state space its training data never
covered. On-policy distillation (OPD) addresses this by training the student on its
own rollouts and having the teacher supply correction on the states the student
actually reaches \autocite{lu_thinkingmachines_opd_2025}. This makes the supervision
signal dense and situated, and OPD has consequently been adopted as a distinct stage
of post-training rather than as a variant of supervised fine-tuning
\autocite{xu_mopd_notion_2025}.

Standard OPD carries two assumptions that are difficult to satisfy in exactly the
case that motivates the problem. The first is white-box access. Per-token correction
requires the teacher's logits, which the strongest available teachers do not expose,
and comparing distributions across model families additionally requires reconciling
different tokenizers, a problem with its own substantial literature and its own
irreducible sources of error \autocite{boizard_uld_2025,minixhofer_alm_2025,phan_tokenization_bias_2025}.
The second assumption concerns granularity. OPD is standardly applied as one
supervision signal over a complete generation, which is adequate for a single
continuous chain of thought but poorly matched to a trajectory composed of discrete,
causally dependent steps.

Existing work relaxes these assumptions separately. Rubric-based On-Policy
Distillation obtains supervision from teacher text alone, using an induced rubric and
a blind verifier in place of logit comparison, and reports strong cross-architecture
results without any tokenizer alignment \autocite{fang_ropd_2026}. Its reward is
nevertheless a single scalar computed once the response is complete. On-Policy
Agentic Distillation, introduced as a component of MAD-OPD, supplies supervision at
the level of individual agentic steps and demonstrates that doing so stabilises
training against error compounding, but obtains that supervision from teacher logits
under a shared vocabulary \autocite{wang_madopd_2026}. OmniOPD operates below
trajectory granularity without logit access, but defines its unit of supervision as
an entropy-selected token span inside a continuous chain of thought, and is evaluated
on single-turn reasoning alone \autocite{zhou_omniopd_2026}.

The combination that the deployment case requires, namely black-box teacher access
together with step-level credit assignment over multi-turn tool-use trajectories, has
not been demonstrated. A practitioner distilling from a proprietary teacher into a
small agent currently chooses between supervision that is dense but unavailable and
supervision that is available but sparse. This thesis addresses that combination.

\section{Problem Statement}
\label{sec:Problem Statement}

Let $\pi_\theta$ denote a student policy and let $\mathcal{T}$ denote a teacher
accessible only as a sampling oracle: given a prompt, $\mathcal{T}$ returns generated
text, and neither its logits, its vocabulary, nor its tokenizer are observable. An
agentic rollout from $\pi_\theta$ produces a trajectory
$\tau = (s_0, a_0, o_0, s_1, a_1, o_1, \dots, s_T, a_T)$, in which $a_t$ is a
model-generated action, typically a tool call together with its accompanying
reasoning, and $o_t$ is the observation returned by the environment. Trajectory-level
black-box supervision assigns a single scalar $r(\tau)$ once the episode terminates.
The credit-assignment problem is that this scalar is uninformative about which $a_t$
was responsible for the outcome, and the number of candidate explanations grows with
$T$.

This thesis asks whether the rubric-based, blind-verifier mechanism that produces
$r(\tau)$ can instead be applied per step to produce a sequence $r(a_t \mid s_t)$
under the same black-box access constraint, and whether training on that denser signal
improves agentic task performance over the trajectory-level baseline. Three questions
follow. First, what constitutes an appropriate unit of supervision in an agentic
trajectory, and how should rubric criteria be induced for a partial rollout whose
eventual outcome is unknown at the time of scoring. Second, whether per-step scoring
introduces failure modes absent at trajectory level, in particular whether rubric
gaming, which the trajectory-level method observes in fewer than 2\% of early-training
rollouts, becomes more prevalent when criteria are local. Third, what the method costs,
given that trajectory-level rubric induction already requires two teacher calls per
rollout and per-step induction multiplies that cost by trajectory length unless the
induction procedure is amortised.

\section{Objectives and Scope}
\label{sec:Objectives}

The objectives of this project are as follows.

\begin{enumerate}[itemsep=2pt,topsep=4pt]
  \item Formalise step-level, black-box supervision for multi-turn agentic
    trajectories, including a defensible definition of the step boundary and a
    procedure for inducing per-step rubric criteria from teacher text.
  \item Implement the resulting method, integrating the per-step reward signal into
    a GRPO training loop.
  \item Evaluate against two baselines on agentic tool-use benchmarks: supervised
    fine-tuning on teacher trajectories, and trajectory-level black-box supervision
    in the manner of ROPD.
  \item Characterise the cost of the method in teacher calls per training step, and
    quantify the trade-off between supervision density and teacher query budget.
  \item Analyse failure modes introduced by per-step scoring, with particular
    attention to rubric gaming and to the effect of scoring steps whose contribution
    to the outcome is not yet determined.
\end{enumerate}

Three directions lie outside the scope of this project. The first is multi-teacher
debate. Translating the ceiling-breaking mechanism of MAD-OPD into black-box space is
a natural continuation of this work and is discussed in
Chapter~\ref{chap:Conclusion}, but a single-teacher method must be established before
a multi-teacher one can be evaluated against it. The second is cross-tokenizer logit
alignment. The method proposed here requires no tokenizer correspondence, and the
cross-tokenizer literature is surveyed in Chapter~\ref{chap:Preliminaries} as context
rather than as a component. The third is the training of new base models. All
experiments fine-tune existing open-weight student models.

\section{Organisation of this Report}
\label{sec:Organisation}

Chapter~\ref{chap:Preliminaries} defines the terminology used throughout, surveys the
relevant literature, and states the contributions of this work.
Chapter~\ref{chap:Implementation} describes the design of the proposed method and its
implementation. Chapter~\ref{chap:Evaluation} presents the experimental setup and
results. Chapter~\ref{chap:Challenges} discusses the difficulties encountered during
the project. Chapter~\ref{chap:Conclusion} summarises the findings and identifies
directions for future work.

